\section{海洋法}

国际法内容最多的一章。

\begin{enumerate}
  \item 领土制度:
        \begin{enumerate}
          \item 构成：陆水（河流、湖泊、三大海域：内海、领海、群岛水域）空底土。
          \item 陆 \yjt 边境制度。
          \item 水 \yjt 河流制度。
          \item 空 \yjt 完全主权，\zd{劫机犯或引渡或起诉}。
          \item 底土 \yjt 完全主权。自然资源。
        \end{enumerate}
        \begin{kk}{取得方式}
          \begin{itemize}
            \item 先占（一定是国家的行为）。天然形成的无主地。\zd{\bdy 最先发现，最先开发}。解决历史问题。
            \item \zd{时效。我国不承认其合法}。
            \item 添附。自然、人工。\zd{以领土为基础}，不能侵害相邻权。
            \item \zd{武力政府违法}。
            \item 割让（看个案）。条约自愿平等与否。\zd{强制不合法}。
          \end{itemize}
          \tcblower
          新的两种：殖民地独立；公民投票（看个案）。国内法有无协议，相关国能否达成协议。
        \end{kk}

  \item 海洋制度：一定会考。海洋水域。海洋底土。放图。

        内水 \yjt 内海 \yjt 领海 \yjt 毗连区 \yjt 专属经济区 \yjt 大陆架 \yjt 公海（国际海域）。

        \begin{enumerate}
          \item \zd{内海完全主权}，非请勿入。
          \item 港口刑事管辖三主动：\zd{扰乱港口安宁，案情重大，受害本国人}，
          \item 领海无害通过权。不允许军舰，连续通过，潜水器浮出水面挂出国旗。
          \item 毗连区：\zd{海关（走私）、财政（逃税）、移民（偷渡）、卫生（检验检疫）、国家安全（我国）}。
          \item 专属经济区：自然资源的专属勘探开发和管辖。事先报备。宣告有效。
          \item 公海：
                \begin{enumerate}
                  \item 船旗国，悬挂多国国旗？
                  \item 普遍（\zd{海盗、非法广播、贩毒、奴隶}）。
                  \item 登临权（起点公海）、紧追权（其他延伸到公海）。起点不同，对象只能非政府、国家船舶。
                  \item 紧追前须发出停止信号（\zd{敲锣打鼓}）、不能中断（接班？）、\zd{追到他国领海}。
                \end{enumerate}
          \item 海洋底土：内海和领海下覆底士、专属经济区下覆底土、大陆架（非生物资源）。
        \end{enumerate}
        \begin{kk}{大陆架}
          \begin{itemize}
            \item 宽度 200 到 350
            \item 举证获得。
            \item 自愿开发内外有别，200 海里以外\zd{国际海底管理局}，\zd{发展中国家}某些条件下享有免缴权。
            \item 平行开发：勘探者和国际海底管理局。
          \end{itemize}
        \end{kk}
        \begin{kk}{群岛水域}
          \begin{itemize}
            \item 岛屿封闭线以外，群岛基线以内，内水和领海之间。
            \item 领水 = 内水 + 领海，完全主权，
            \item 无害通过制度。船无害通过，飞机经许可，空中通道（权利）不经许可。
          \end{itemize}
        \end{kk}
        \begin{kk}{国际海峡}
          \begin{enumerate}
            \item 两端连接的是公海 \hz 专属经济区。
            \item 通行制度（最最重要）。
                  \begin{enumerate}
                    \item 过境通行制：飞机船舶不经许可不打招呼无害通过。
                    \item 无害通过权（船）：一国大陆和该国岛屿构成的国际海峡且该岛屿向海一面有方便的航道。
                    \item 如果够宽中间航道，自由通行制，走走停停。
                    \item 沿岸国可以\zd{指定航道、规定分道分航}（是权利）。
                  \end{enumerate}
          \end{enumerate}
        \end{kk}
  \item 特殊空间法律制度。
        \begin{kk}{南极}
          可考性相对差一些。
          \begin{enumerate}
            \item 人类共同利益原则、和平目的原则、自由科考原则。
            \item \zd{冻结（不前进不后退）}南极领土主权要求。
          \end{enumerate}
        \end{kk}
        \begin{kk}{外空}
          航空器能到达的高度。
          \begin{enumerate}
            \item 登记制度（联合国秘书长）。登记产生法律后果：管辖权和所有权的依据。
            \item 营救制度：\zd{援助、通知、送回}。
            \item 责任制度：\zd{发射国：扩大解释}。
                  \begin{enumerate}
                    \item 空对地绝对责任，空对空过错责任。
                    \item 单方过错单方承担，都有过错按比例，没有过错损失自担。
                    \item \zd{不调整缔约国和本国国民，缔约国和发射国}。
                  \end{enumerate}
          \end{enumerate}
        \end{kk}
  \item 国际环保法：
        \begin{enumerate}
          \item 大气环保（重点热点）：
                \begin{enumerate}
                  \item 92 年联合国气候变化框架公约、97 年京都议定书、15 年巴黎协定。
                  \item \zd{共同但有区别责任原则}。
                  \item 工业化国家、发达国家（\zd{拿出资金和技术援助}）、发展中国家。
                  \item 净排量：实际排量扣掉本土森林所能吸收的二氧化碳数量。
                  \item 三种折抵方式：\zd{排放权交易 \yjt 发达国家，输出绿色技术 \yjt 发展中国家,集团方式 \yjt 欧盟}。
                  \item 2020 年以后，国家自主贡献方式，有减排义务，但不再适用公约的强制摊派指标。
                \end{enumerate}
          \item 濒危野生动植物种国际贸易公约。
                \begin{kk}{以清单为基础的许可证制度}
                  \begin{enumerate}
                    \item 附件一：进出口许可证。
                    \item 附加二：出口许可证。
                    \item 附件三：自行决定管理。
                    \item 活体和标本。
                  \end{enumerate}
                \end{kk}
          \item 控制危险废物越境转移（巴塞尔公约）。
                \begin{enumerate}
                  \item \zd{缔约国}之间。
                  \item 特定进口\zd{书面}同意。
                  \item 有无害环境的处置方法（核废水）。
                  \item 必须要\zd{有保险 \hz 担保}。
                \end{enumerate}
        \end{enumerate}
\end{enumerate}
